% Engineering observations. Included in the results (journal) or the appendix (conference).
\paragraph{Policy-compatible rationale fields} With the runtime's field names exposed as they were, every task
of an initial mock run was declined at its first request with the refusal category
\texttt{reasoning\_extraction}. After presenting the fields as \texttt{purpose} and \texttt{summary} and removing the instruction to
state hypotheses from the first message (\S\ref{sec:alias}), all tasks completed, and the runtime still
received its required \texttt{reasoning} fields.

\paragraph{Parallel execution} On a simulated runtime that allowed two tasks in flight, the agent was started
with four workers. It learned the limit from the first refused leases (the log shows the limit lowered first
to three, then to two, as refusals arrived), never ran more than two tasks at once afterwards, and completed
all 16 tasks of that run in 112\,s with the real model, without a failed task.

\paragraph{Benign twins} The two Linux service tasks share an identical alert (``Service state changed by
systemctl as root'') and graph, and differ only in the lineage of the triggering process: an operator's login
shell through \texttt{sudo} versus the package manager's post-installation script. They test whether an agent
reads provenance rather than the alert text. The agent never called the package manager's activity malicious,
but it cleared the operator's service stop in 7 of 17 runs: an interactive administrator stopping a web server
is an ordinary event in isolation, and only the lab context makes it an attack. Twin pairs thus expose the
direction in which an agent errs when provenance alone is ambiguous.
